\begin{abstract}
Solving novel view synthesis (NVS) for millimeter-wave (mmWave) radar requires a renderer that is multi-viewpoint-tractable, physically faithful, and complex-valued. No prior method achieves these three properties simultaneously. Physically based radar renderers must condition radar synthesis not only on sensor pose, but also on the full antenna array geometry, waveform configuration, and beam patterns. Differentiable Monte Carlo (MC) ray tracers implement the radar forward model directly with explicit material modeling and complex outputs, but do not scale to the multi-view optimization NVS demands. Optical-NVS ports of NeRF, hash grids, and 3D Gaussians train fast but discard phase and replace explicit material modeling with opaque learned features, restricting them to power-only range--azimuth (RA) magnitudes. We propose 3D Point Splatting (3DPS), the first differentiable point renderer for radar, derived directly from the standard solid-angle form of the radar equation. Each oriented 3D point carries an ITU-R~P.2040 material model, evaluated in closed form, with the resulting complex phasor splatted into range bins through a precomputed point spread function (PSF). Complex-valued output makes the renderer product-agnostic: the same optimized scene yields ADC, complex range profile (CRP), and RA outputs through standard FFT pipelines without retraining for each format. On six outdoor ColoRadar scenes, 3DPS reaches \textbf{0.587} mean Pearson correlation on held-out RA images, $1.8$--$7.2\times$ over RadarSplat, Radar Fields, and DART. Training takes ${\sim}3$~minutes per scene on a single RTX~4090, ${\sim}50\times$ faster than the SOTA MC ray tracer's projected per-scene cost.
\end{abstract}
